\documentclass[letterpaper]{article}
% AAAI 2027 style is not present locally. Swap this preamble to the official
% AAAI 2027 (or AAAI 2026 stand-in) style file before submission packaging.
\usepackage[utf8]{inputenc}
\usepackage[T1]{fontenc}
\usepackage{times}
\usepackage{helvet}
\usepackage{courier}
\usepackage{url}
\usepackage{graphicx}
\usepackage{booktabs}
\usepackage{amsmath}
\usepackage{amssymb}
\usepackage{multirow}
\usepackage{xspace}

\title{PhysMon: Monitoring Shortcut Sensitivity in Scientific Reasoning from Prompt-Side Hidden States}
\author{[Authors Redacted for Drafting]}
\date{}

\begin{document}
\maketitle

\begin{abstract}
We study whether prompt-side hidden states of open-weight language models can predict model-specific shortcut sensitivity on physics problems before any failure appears in generated output. We construct PhysMon, a family-structured benchmark of 140 canonical counterfactual families in classical mechanics, electrostatics, and thermodynamics, where designated cue edits are certified invariant by a symbolic solver. Qwen2.5-7B-Instruct shows non-trivial output instability on 50.0\% of families (mean $S^{\mathrm{lp}}=1.4006$ nats); Llama-3.1-8B on 22.9\%; DeepSeek-R1-Distill-Qwen-32B on 27.1\%. A variance probe trained on prompt-side hidden states at layer 18 predicts family-level sensitivity with AUROC 0.731 (ensemble 0.748), outperforming surface and textual baselines including CoT classifiers (0.5990), answer+rationale classifiers (0.6464), and LLM judges (zero-shot AUROC 0.439--0.500 across settings). Controlled multi-head knockout of heads \{26, 24, 13, 11\} at exactly layer 16 completely suppresses sensitivity in 17/20 tested families, with causal specificity confirmed against random-direction, layer, and construct controls. Sensitivity encoding localises at 64--69\% network depth consistently across Qwen-7B and DeepSeek-32B despite a 4.5$\times$ scale difference and different training regimes. Reasoning training selectively suppresses unit-compatible distractor sensitivity (Fisher $p=0.00575$, OR $=\infty$) while leaving frame-rendering sensitivity intact. Correctness-residualised AUROC is [TBD-RESID]. We release benchmark templates, solver derivations, and analysis code.
\end{abstract}

\section{Introduction}
% Motivation: shortcut sensitivity in scientific reasoning; why output-only evaluation misses it.
% Introduce PhysMon, prompt-side monitoring, and the main empirical story.
% Finalize with paper contributions list.

\section{Background and Related Work}
% Hidden-state monitoring and probes.
% Scientific reasoning benchmarks and counterfactual evaluation.
% Mechanistic interpretability and causal interventions.
% Reasoning-tuned models and chain-of-thought robustness.

\section{Formal Framework}
% Define family-structured counterfactual benchmark.
% Three-level framework: prompt condition, observed behaviour, internal representation.
% Formal notation for S, hat{S}, and S^{lp}.
% Note where H1--H4 are tested and what remains scoped as conditional.

\section{PhysMon Benchmark}
% Benchmark construction, domains, cue types A/B/C, and symbolic verification.
% Family counts: 140 total; mechanics/electrostatics/thermodynamics split.
% Near-match distractor principle and Cue C framing.
% Mention family validation pipeline and remaining human-audit follow-up.

\section{Hidden-State Monitoring Method}
% Activation extraction sites and why last-prompt residual variance is primary.
% Variance probe, ensemble probe, domain generalisation, cross-model transfer.
% Baselines: TF-IDF, CoT text classifier, answer+rationale classifier, correctness probe,
% LLM judges, self-consistency (pending final integration from job 244037).
% Residualisation analysis placeholder: add when job 245081 completes.

\section{Results}

\subsection{Behavioural Sensitivity}
% Report model-level positive rates and mean S^{lp}.
% Highlight near-match distractor effect and Cue C separation.
% DeepSeek selective suppression/amplification story.

\subsection{Hidden-State Monitoring}
% Primary AUROC 0.731, ensemble 0.748, Pearson r = 0.667.
% No-Cue-C result 0.715 with delta +0.104 over no-Cue-C TF-IDF.
% Domain generalisation macro AUROC 0.749.
% Correctness-probe anti-correlation analysis and residualisation [TBD-RESID].

\subsection{Causal Representation}
% Full-residual patching, H26 knockout, and 4-head L16 knockout.
% Random-direction, unrelated-donor, early-layer, damage-control, and injection results.
% Phrase causal claim carefully: necessity strongly supported; sufficiency directional but incomplete.

\subsection{Cross-Architecture and Reasoning Training}
% DeepSeek within-model variance probe AUROC 0.7025.
% Depth universality: Qwen 18/28 vs DeepSeek 44/64.
% Qwen->DeepSeek transfer 0.7252 vs Qwen->Llama 0.637.
% Latent-vs-novel sensitivity split in DeepSeek-amplified families.

\section{Discussion}
% Interpret what the probe is detecting mechanistically.
% Distinguish shortcut sensitivity from generic difficulty.
% Discuss reasoning training as selectively suppressive rather than uniformly robustifying.
% State limitations: residualisation pending, self-consistency pending, same-answer/no-cue donor deferred,
% variable-renaming/human-audit still to be completed before submission freeze.

\section{Conclusion}
% Restate main finding and honest scope of the current claims.
% Emphasize prompt-side monitorability and circuit-level necessity result.

% Figures/tables to insert:
% - Baselines table: paper/tables/baselines_table.tex
% - Causal controls table: paper/tables/causal_controls_table.tex
% - Layer AUROC curves, MHK heatmap, injection dose-response, DeepSeek cue-type plot.

\bibliographystyle{plain}
\bibliography{references}

\end{document}
